% KHÔNG \input file này vào main.tex / chuong2.tex.
% Toàn bộ mục ``X là gì'' kiểu bách khoa dễ trùng Turnitin với Wikipedia,
% tài liệu hãng và khóa luận CNTT khác. Chương 2 đã giải thích bằng việc đồ án
% cần làm; không cần thêm lớp định nghĩa này.

\section{Định nghĩa các công nghệ và lý thuyết cốt lõi}

Các mục trước đã gắn kỹ thuật vào bài toán pháp lý. Mục này trả lời đúng dạng
câu hỏi ``X là gì'': khái niệm, đặc điểm, rồi mới nói đồ án dùng X chỗ nào.
Cách viết này phục vụ người đọc chưa làm RAG hoặc chưa dựng cụm Docker, đồng
thời giữ giọng báo cáo (không phải hướng dẫn cài đặt từng lệnh).

\subsection{RAG là gì}

Retrieval-Augmented Generation (RAG) là kiến trúc kết hợp hai khối: (1)~truy hồi
tài liệu liên quan từ một kho bên ngoài mô hình; (2)~sinh câu trả lời với tài
liệu đó nằm trong ngữ cảnh~\cite{lewis2020rag}. Khác chatbot close-book, RAG
không lấy ``trí nhớ'' nén trong trọng số làm nguồn chính. Khi tri thức đổi (luật
sửa, thông tư mới), người vận hành cập nhật kho, không bắt buộc huấn luyện lại
mạng.

Ba thành phần tối thiểu: kho đã cắt đoạn; bộ truy hồi (thưa, dày, hoặc lai);
mô hình sinh có cửa sổ ngữ cảnh đủ chứa đoạn đã chọn. RAG không tự bảo đảm đúng
luật: nếu truy hồi sai, câu sinh ra có thể trôi mà vẫn sai. Đồ án thêm hậu kiểm
số Điều so với tập vừa tìm, và lọc LLM trước khi sinh.

Đồ án hiện thực RAG trên đơn vị Điều, hybrid BM25--Chroma, RRF, HyDE, bảy nút
LangGraph. Đây là RAG miền hẹp (closed corpus), không phải RAG Internet.

\subsection{Mô hình ngôn ngữ lớn (LLM) là gì}

LLM là mạng (thường Transformer decoder) huấn luyện trên khối văn bản rất lớn,
nhiệm vụ gốc là dự đoán token tiếp theo. Từ đó mô hình viết được câu, tóm tắt,
dịch, trả lời hỏi. LLM không có cơ chế ``mở Điều 44'' trừ khi đoạn Điều được
đưa vào prompt hoặc đã xuất hiện dày trong dữ liệu huấn luyện --- trường hợp
sau không kiểm chứng được và dễ bịa số hiệu.

Đồ án gọi LLM qua API (Gemini Flash), nhiệt độ 0, không fine-tune. LLM đóng vai
viết câu, phân nhãn intent/PASS-DROP, và viết đoạn HyDE; không đóng vai CSDL
luật.

\subsection{Cơ sở dữ liệu vector là gì}

CSDL vector lưu embedding và tìm $k$ láng giềng gần theo cosine hoặc L2. Khác
CSDL quan hệ: không tối ưu JOIN, ràng buộc khóa ngoại, hay giao dịch kế toán.
Đồ án dùng Chroma cho vector, Postgres cho nghiệp vụ. Một hệ không thay hệ kia.

\subsection{ReactJS là gì}

React (thường gọi ReactJS khi nhấn môi trường JavaScript trình duyệt) là thư
viện UI do Meta phát triển~\cite{react}. Đơn vị xây dựng là component: hàm nhận
dữ liệu (props) và state, trả về mô tả giao diện (JSX). JSX trông giống HTML
trong JavaScript nhưng được biên dịch thành lời gọi
\texttt{React.createElement}. Khi state đổi, React tính lại cây và cập nhật DOM
thật theo diff (reconciliation), tránh thao tác DOM thủ công dễ lệch.

React không phải khung MVC đầy đủ: không có router, không có HTTP client, không
có CSDL. Hệ sinh thái bổ sung (React Router, TanStack Query, Zustand) đi kèm.
Ứng dụng React điển hình là SPA: một lần tải JS, đổi trang không reload.

Đồ án dùng React 19 cho toàn bộ màn hình sau đăng nhập. Không dùng Next.js vì
không cần SSR. Component tách theo nghiệp vụ (chat, citation, cây luật) để test
parser SSE độc lập.

\subsection{JSX, props, state}

Props là dữ liệu cha truyền xuống, coi như đầu vào hàm, không sửa tại chỗ.
State là dữ liệu component tự giữ, đổi bằng \texttt{setState}/\texttt{useState},
kéo theo render. Lẫn hai khái niệm làm dữ liệu một chiều bị phá (con sửa props
của cha). Đồ án: danh sách citation là props từ trang chat; ô nhập là state
cục bộ; user đăng nhập là state Zustand toàn cục.

\subsection{Docker là gì}

Docker đóng gói ứng dụng cùng thư viện, runtime và cấu hình thành image, chạy
thành container trên kernel máy chủ~\cite{docker}. Khác máy ảo: không mang cả
hệ điều hành khách đầy đủ, nên khởi động nhanh và tốn ít RAM hơn. Dockerfile
mô tả các lớp (cài Python, copy mã, lệnh chạy). Lớp không đổi được cache, build
lần sau nhanh hơn.

Container mặc định mất dữ liệu khi xóa; volume gắn thư mục bền. Mạng Docker cho
các container gọi nhau bằng tên dịch vụ (\texttt{postgres}, \texttt{backend}).

Đồ án: một image backend Python 3.10, một image frontend nginx, một image
Postgres chính thức. Người chấm máy khác chạy cùng phiên bản phụ thuộc, không
cần cài Node/Python trên host.

\subsection{Docker Compose là gì}

Compose là công cụ mô tả \emph{cụm} nhiều container bằng YAML: dịch vụ, cổng,
biến môi trường, volume, phụ thuộc, healthcheck. Một lệnh \texttt{up} tạo mạng
và khởi động đúng thứ tự (Postgres healthy rồi mới backend). Đồ án ba dịch vụ;
Compose đọc \texttt{.env} ở thư mục gốc dự án.

\subsection{FastAPI là gì}

FastAPI là framework web Python hiện đại, dựa ASGI, tự sinh OpenAPI, kiểm tra
dữ liệu bằng Pydantic~\cite{fastapi}. Phù hợp API JSON và luồng bất đồng bộ.
Không nhằm render HTML server-side (khác Django template). Đồ án dùng FastAPI
làm cửa duy nhất cho SPA: REST + SSE, không phục vụ JSX.

\subsection{PostgreSQL là gì}

PostgreSQL là hệ quản trị CSDL quan hệ nguồn mở: SQL, giao dịch ACID, khóa
ngoại, chỉ mục, kiểu mở rộng (JSONB, tsvector). Đồ án bản 17. Không dùng như
``kho file luật thô'': mỗi Điều là một hàng có ràng buộc và FTS. Mã số thuế
unique, task tuân thủ unique theo kỳ --- việc này CSDL bắt, không chỉ tin Python.

\subsection{LangGraph là gì}

LangGraph mô hình hóa tác nhân thành đồ thị: nút là bước xử lý, cạnh là chuyển
trạng thái, state đi qua các nút. Hợp pipeline có nhánh (SKIP retrieval) và cần
nhật ký từng bước. Khác ``một hàm 400 dòng'': mỗi nút test và quan sát được.
Đồ án bảy nút cố định thứ tự, không phải agent tự chọn công cụ tự do.

\subsection{TypeScript là gì}

TypeScript là JavaScript có hệ thống kiểu tĩnh, biên dịch ra JS chạy trên trình
duyệt. Kiểu giúp bắt lệch tên trường API lúc build. Đồ án khai báo interface
cho user, Điều, finding, event SSE. Không chạy TypeScript trên server sản phẩm
frontend --- nginx chỉ phục vụ JS đã dịch.

\subsection{JWT là gì}

JSON Web Token là chuỗi ba phần (header, payload, chữ ký) dùng mang thông tin
xác thực giữa hai bên~\cite{rfc7519}. Máy chủ ký payload; lần sau chỉ kiểm chữ
ký và hạn, không cần bảng session trên RAM. Đồ án HS256, hai loại token
(access/refresh) phân biệt claim \texttt{type}. JWT không mã hoá payload mặc
định: không nhét mật khẩu vào token.

\subsection{REST là gì}

Representational State Transfer là phong cách thiết kế API: tài nguyên có URL,
thao tác bằng HTTP method, không trạng thái server giữa hai request. Đồ án:
\texttt{/api/v1/laws} là tài nguyên văn bản; GET đọc, POST tạo. REST không cấm
SSE: hỏi đáp là POST tạo lượt hỏi, đáp trả dạng stream.

\subsection{nginx là gì}

nginx là máy chủ HTTP, thường dùng reverse proxy và phục vụ tệp tĩnh. Đồ án:
image frontend là nginx đọc bản Vite build, chuyển \texttt{/api} sang container
backend, tắt buffering cho SSE. Người dùng chỉ mở một cổng 5173 trên host.

\subsection{Vite là gì}

Vite là công cụ dựng frontend: dev server ESM nhanh, production bundle bằng
Rollup. Đồ án không tự viết webpack. Biến \texttt{VITE\_*} được nhúng lúc build,
không lúc chạy nginx --- để trống base URL trong Docker.

\subsection{Tailwind CSS là gì}

Tailwind là framework CSS kiểu utility: class nhỏ (\texttt{flex},
\texttt{px-4}) gắn trong JSX. Khác Bootstrap (component có sẵn). Đồ án tự bố
cục, bản 4, JIT chỉ phát CSS dùng tới.

\subsection{ChromaDB là gì}

Chroma là kho vector nhúng được, API thêm/tìm embedding kèm metadata. Đồ án
chạy trong process backend, dữ liệu volume \texttt{chroma\_db}. Không thay
Postgres. Manifest mô hình/chiều phải khớp lúc khởi động.

\subsection{BM25 là gì}

BM25 là hàm xếp hạng xác suất trong truy hồi thông tin: từ nào hiếm (IDF cao)
và xuất hiện vừa đủ trong tài liệu thì điểm cao, tài liệu dài bị phạt~\cite{robertson2009bm25}.
Không cần GPU. Yếu khi từ vựng lệch. Đồ án dùng BM25 song song vector.

\subsection{HyDE là gì}

HyDE (Hypothetical Document Embeddings) nhờ LLM viết một đoạn ``như thể đã trả
lời'', rồi nhúng đoạn đó để tìm vector~\cite{gao2023hyde}. Đoạn giả có thể sai
nội dung nhưng gần giọng kho. Đồ án bật HyDE, không hiện đoạn giả cho người dùng
như đáp án.

\subsection{RRF là gì}

Reciprocal Rank Fusion gộp nhiều bảng xếp hạng bằng tổng nghịch đảo hạng, không
cần chuẩn hoá điểm~\cite{cormack2009rrf}. Đồ án $k=60$, trọng số dense:BM25
$=3:1$.

\subsection{SSE là gì}

Server-Sent Events: kết nối HTTP máy chủ chủ động đẩy sự kiện văn bản
(\texttt{event:}, \texttt{data:}). Một chiều. Đồ án: \texttt{status},
\texttt{token}, \texttt{result}. Client đọc Fetch vì EventSource không gắn Bearer.

\subsection{Argon2 là gì}

Argon2 là hàm băm mật khẩu thắng Password Hashing Competition, tốn bộ nhớ và
thời gian có chủ ý để chống dò GPU~\cite{argon2}. Đồ án argon2id qua passlib.
Không dùng SHA-256 cho mật khẩu.

\subsection{ORM là gì}

Object-Relational Mapping ánh xạ bảng thành lớp. CRUD thành thao tác đối tượng.
Đồ án SQLAlchemy 2 async. ORM không miễn viết hiểu schema: generated column,
GIN, unique task vẫn phải thiết kế đúng.

\subsection{Alembic là gì}

Alembic ghi lại thay đổi schema thành revision có mã cha, máy khác chạy
\texttt{upgrade} ra cùng cột. Đồ án không sửa tay bảng trên container rồi quên
ghi.

\subsection{Pydantic là gì}

Thư viện khai báo mô hình dữ liệu bằng type hint, kiểm tra lúc chạy, xuất JSON.
Là lớp biên giới API đồ án: body sai thì 422, không vào nghiệp vụ.

\subsection{OpenAPI là gì}

Chuẩn mô tả API (đường, schema, mã lỗi) để máy và người đọc. FastAPI sinh từ
hint. Đường \texttt{/docs} dùng khi bảo vệ, không phải giao diện user DNNVV.

\subsection{SPA là gì}

Single Page Application tải một bộ HTML/JS, điều hướng client, gọi API lấy dữ
liệu. Ưu: UI mượt. Nhược: SEO kém --- chấp nhận vì trang sau login. Đồ án SPA
React + nginx fallback \texttt{index.html}.

\subsection{CORS là gì}

Cơ chế trình duyệt kiểm soát đọc đáp khác origin. Server gửi
\texttt{Access-Control-Allow-Origin}. Đồ án danh sách origin, không \texttt{*}.
Docker cùng origin thì CORS gần như không chạy.

\subsection{Unicode và UTF-8 là gì}

Unicode gán mã điểm cho ký tự mọi ngôn ngữ. UTF-8 mã hoá biến độ dài, tương
thích ASCII. Toàn bộ Điều luật đồ án UTF-8. Tìm không dấu là bản phụ, không phải
đổi encoding hiển thị.

\subsection{ACID là gì}

Atomicity, Consistency, Isolation, Durability --- bốn tính chất giao dịch quan
hệ. Đồ án cần khi đăng ký kèm tạo tổ chức: không để nửa user nửa công ty.

\subsection{Chỉ mục CSDL là gì}

Cấu trúc phụ giúp tìm hàng không quét cả bảng. B-tree cho so khớp khóa; GIN cho
tsvector và trigram. Đồ án lập cả hai nhóm trên bảng Điều.

\subsection{Connection pooling là gì}

Giữ sẵn kết nối TCP tới Postgres, cho mượn theo request. Mở kết nối mới mỗi API
sẽ chậm và cạn \texttt{max\_connections}. SQLAlchemy async dùng pool.

\subsection{XSS và CSRF là gì}

XSS chèn JavaScript vào trang nạn nhân, đọc \texttt{localStorage}. CSRF mượn
cookie phiên gửi request giả. Đồ án: React escape text; Bearer header giảm CSRF
so với cookie cổ điển; XSS vẫn là lý do ghi nợ chuyển cookie HttpOnly.

\subsection{Rate limiting là gì}

Giới hạn số request theo thời gian. Bảo vệ máy chủ và hạn mức LLM. Đồ án
slowapi: chat 20/phút, upload 30/giờ.

\subsection{Healthcheck là gì}

Lệnh định kỳ (Docker/Compose) hỏi dịch vụ còn sống không. Đồ án:
\texttt{pg\_isready} và HTTP \texttt{/health} không gọi Gemini. Backend đợi
Postgres healthy trước khi phục vụ.

\subsection{Reverse proxy là gì}

Máy chủ đứng trước, nhận HTTP từ client, chuyển tiếp server nội bộ. Che cổng
backend, gộp origin, có thể nén và TLS. Đồ án: nginx.

\subsection{Background task là gì}

Việc chạy sau khi đã trả HTTP, không giữ trình duyệt chờ. FastAPI
\texttt{BackgroundTasks} cùng tiến trình --- đủ đồ án, mất việc nếu container
chết giữa soát xét. Khác hàng đợi Redis/Celery.

\subsection{Prompt engineering là gì}

Soạn lệnh hệ thống và người dùng để mô hình bám nhiệm vụ, không đổi trọng số.
Đồ án: bốn vai prompt (intent, HyDE, filter, generate), đầu ra hẹp
(NEXT/SKIP, PASS/DROP). Fine-tune không dùng.

\subsection{Fine-tuning là gì}

Cập nhật trọng số mô hình trên dữ liệu miền. Tốn GPU; luật đổi phải train lại.
Đồ án thay bằng nhập corpus + prompt.

\subsection{Token và tách từ là gì}

Token là đơn vị mô hình hoặc BM25 đếm. Tiếng Việt từ không trùng khoảng trắng;
underthesea tách từ cho BM25. LLM dùng subword của nhà cung cấp --- hai lớp
token khác nhau, đồ án chấp nhận.

\subsection{Hallucination là gì}

Mô hình sinh nội dung không có trong ngữ cảnh hoặc không có thật. Pháp lý: bịa
số Điều, bịa mức phạt. RAG giảm nhưng không triệt. Hậu kiểm citation là lớp
xác định bù lớp xác suất.

\subsection{Open-domain QA là gì}

Hỏi đáp không gắn sẵn một đoạn văn: hệ thống phải tự tìm đoạn trong kho lớn rồi
đọc~\cite{karpukhin2020dpr}. Đồ án là QA miền hẹp (38 văn bản) nhưng cùng kiến
trúc retrieve-then-read.

\subsection{F2 là gì}

$F_\beta$ với $\beta=2$ nhấn độ phủ hơn độ chính xác. Hướng retrieval pháp lý:
thà lấy thừa Điều (lọc sau) hơn bỏ sót Điều quyết định. Đồ án dùng F2 như khái
niệm thiết kế, chưa công bố điểm F2 vì chưa gắn nhãn bộ câu nội bộ đủ lớn.

\subsection{uv là gì}

Công cụ quản lý gói và môi trường ảo Python (tương tự pip+venv, khóa phụ thuộc
chặt). Đồ án dùng uv trên máy dev và trong Dockerfile backend, thống nhất
\texttt{pytest}.

\subsection{Node.js và npm là gì}

Node.js chạy JavaScript phía máy phát triển (build, test). npm cài gói
\texttt{package.json}. Sản phẩm Docker frontend \emph{không} chạy Node: chỉ
nginx và tệp tĩnh. Nhầm ``React cần Node trên server'' là sai trong kiến trúc
này.

\subsection{Git là gì}

Hệ thống quản lý phiên bản phân tán: commit, nhánh, diff. Dùng khi phát triển
đồ án để hoàn tác và đối chiếu schema. Báo cáo không phụ thuộc remote công khai.

\subsection{WSL2 là gì}

Windows Subsystem for Linux 2: máy ảo Linux nhẹ trên Windows, dùng làm host
Docker. Không phải thành phần sản phẩm. Linux thuần chạy cùng Compose.

\subsection{ASGI là gì}

Giao diện bất đồng bộ giữa máy chủ Python và ứng dụng, hỗ trợ HTTP dài và
WebSocket/SSE. FastAPI chạy ASGI (Uvicorn). WSGI cổ điển không đủ cho SSE đồ án.

\subsection{Uvicorn là gì}

Máy chủ ASGI chạy ứng dụng FastAPI. Docker backend: Uvicorn không reload, cổng
8023. Số worker giữ một tiến trình vì Chroma nhúng in-process.

\subsection{asyncpg là gì}

Driver PostgreSQL bất đồng bộ cho Python. SQLAlchemy 2 dùng URL
\texttt{postgresql+asyncpg://}. Không chặn vòng sự kiện khi chờ SQL.

\subsection{passlib là gì}

Thư viện băm và kiểm mật khẩu, bọc argon2. \texttt{verify} trả False khi hash
hỏng, không làm 500 trang login.

\subsection{Matryoshka embedding là gì}

Kỹ thuật huấn luyện để tiền tố vector vẫn mang nghĩa, cho phép cắt số chiều.
Đồ án cắt 1536 từ 3072 chiều Gemini. Không phải mọi mô hình cắt được.

\subsection{Backoff là gì}

Khi API trả 429 hoặc 5xx, chờ dần (exponential) rồi thử lại, có trần số lần.
Đồ án dùng khi nhúng lô Điều. Không thử vô hạn kẻo treo container.

\subsection{Idempotency là gì}

Gọi lại cùng thao tác không tạo thêm hiệu ứng ngoài lần đầu. Sinh lịch đồ án
idempotent theo (tổ chức, rule, kỳ). Gửi chat không idempotent: mỗi POST một
tin mới.

\subsection{multipart/form-data là gì}

Kiểu HTTP tải file: boundary tách nhị phân và trường form. Axios phải để trình
duyệt tự gắn Content-Type. Đồ án interceptor xóa header JSON khi body là
\texttt{FormData}.

\subsection{OpenAI-compatible API là gì}

Nhiều thư viện nói chuyện theo đường \texttt{/chat/completions} và
\texttt{/embeddings} kiểu OpenAI. Gemini cung cấp lớp tương thích; đồ án đổi
\texttt{base\_url} là gọi được, sau này trỏ vLLM không viết lại node.
